% =====================================================================
% paper/sections/theoretical_model.tex
% Theoretical Model: Two-Country OLG with Demographics and Public Debt
% Self-contained section -- to be \input{} from paper/main.tex
% Preamble (working-paper-format.md) assumed: amsmath, amsthm,
% amssymb, mathtools, dsfont, biblatex+biber, cleveref, booktabs,
% tabularray, threeparttable.
% =====================================================================

\section{Theoretical Model}\label{sec:theory}

\subsection{Overview and Logic of the Model}\label{sec:theory:overview}

We develop a two-country overlapping-generations (OLG) model in the tradition of
\citet{Samuelson1958} and \citet{Diamond1965}, augmented with stochastic
demographic structure, social-security transfers, public debt, and
convex-adjustment-cost investment \citep{Hayashi1982}. The model's purpose is
not to characterise dynamics or welfare but to deliver \emph{signed
predictions} for the panel coefficients of a long-run regression linking
each country's net international investment position (NIIP) to demographic,
fiscal, and productivity fundamentals. The empirical specification is therefore
a direct image of a structural object derived below.

The logic proceeds in four steps. First, households of each generation solve a
CRRA life-cycle problem (\cref{sec:theory:households}); aggregation over young
and old yields private saving (\cref{sec:theory:agg-saving}). Second, firms
choose investment subject to convex adjustment costs, yielding a Tobin-$q$
representation (\cref{sec:theory:firms}). Third, world capital-market clearing
determines the world real interest rate $r_t^*$, and the implicit-function
theorem (IFT) signs its dependence on global fundamentals
(\cref{sec:theory:clearing}). Fourth, the country-level intertemporal budget
constraint (IBC) for the NIIP is forward-iterated under a no-Ponzi condition;
along the balanced-growth path (BGP) the IBC collapses to a perpetuity that we
linearise around country-specific reference points
(\cref{sec:theory:linearization}). The linearisation produces the regression
equation, with each coefficient given by a partial derivative of structural
saving minus investment, scaled by $(r^*-g_i^Y)^{-1}$. The error-correction
representation (\cref{sec:theory:ecm}) then follows from the Granger
representation theorem \citep{EngleGranger1987}.

\paragraph{Notation conventions.} \Cref{tab:notation} summarises the symbols
used throughout. We adopt the conventions: $\rho_i$ for the country-$i$
error-correction speed (so $\alpha$ is reserved exclusively for the capital
share); $H$ for the expectation horizon (so $T$ can be reserved for sample
length in the empirical section); $b_{it}^{LR}$ for the long-run (BGP) NIIP
(so the asterisk on $r_t^*$ is unambiguous and refers only to the world
interest rate); $g_i^Y \equiv g_i^A + g_i^L$ for the country-$i$ BGP growth
rate of output.

\begin{table}[t]
\centering
\caption{Notation: Symbols and Conventions}\label{tab:notation}
\begin{threeparttable}
\begin{tabular}{ll}
\toprule
Symbol & Meaning \\
\midrule
$i\in\{H,F\}$ & Country index (home / foreign) \\
$t$ & Calendar time index \\
$L_{it}$ & Working-age population, country $i$ at $t$ \\
$\phi_{it}$ & Old-age dependency ratio, $L_{i,t-1}/L_{it}$ \\
$\omega_{it}$ & Effective retirement duration (years old / years working) \\
$A_{it}$, $\tilde a_{it}$ & TFP level and (log) deviation from world mean \\
$g_{it}^A,\ g_{it}^L$ & Growth rates of TFP and population \\
$g_i^Y \equiv g_i^A+g_i^L$ & BGP growth rate of output, country $i$ \\
$\alpha$ & Capital share in Cobb--Douglas production \\
$\delta$ & Depreciation rate \\
$\beta,\ \sigma$ & Discount factor; CRRA intertemporal elasticity \\
$r_t^*$ & World real interest rate (gross: $1+r_t^*$) \\
$w_{it}$ & Competitive real wage \\
$s_{it}^*$ & \emph{Individual} saving of a young agent of generation $t$ \\
$tr_{it}$ & Per-old-agent pension transfer \\
$\tau_{it}$ & Per-young-agent payroll tax \\
$d_{it}^g$ & Public debt per efficiency worker \\
$B_{it}$, $b_{it}$ & NIIP (level); NIIP per efficiency worker \\
$\iota_{it}$, $s_{it}^Y$ & Investment and national saving per efficiency worker \\
$ca_{it}^{NR}$ & Non-interest current account, $s_{it}^Y-\iota_{it}$ \\
$b_{it}^{LR}$ & Long-run (BGP) NIIP, the structural object of interest \\
$\rho_i$ & Country-$i$ error-correction adjustment speed \\
$H$ & Forward expectation horizon (for $\phi,\omega$) \\
$\lambda_t,\mu_i$ & Time and country fixed effects in the panel \\
\bottomrule
\end{tabular}
\begin{tablenotes}\footnotesize
\item \textit{Notes.} The capital share $\alpha$ and the ECM adjustment speed
$\rho_i$ are kept notationally distinct (correcting an overload in earlier
drafts). The expectation horizon $H$ replaces an earlier use of $T$ to avoid
collision with sample length. $b_{it}^{LR}$ replaces $b_{it}^*$ to avoid the
asterisk collision with $r_t^*$.
\end{tablenotes}
\end{threeparttable}
\end{table}

\subsection{Economic Environment}\label{sec:theory:environment}

We list nine assumptions. Eight describe the primitive environment; the ninth
is a stability condition on the world capital market that was implicit in
earlier drafts and that we now name and invoke explicitly in
\cref{prop:cs-rstar}.

\begin{assumption}[Single Tradable Good]\label{ass:A1}
There is a single tradable consumption good, freely traded across countries at
a unit world price. The real exchange rate is therefore indeterminate in the
model; we control for it empirically.
\end{assumption}

\begin{assumption}[Two-Period OLG with Demographic Heterogeneity]\label{ass:A2}
Each country is populated by overlapping generations indexed by birth date.
Agents live two periods (young, old). The size of generation $t$ in country
$i$ is $L_{it}$, so the old-age dependency ratio is $\phi_{it}\equiv L_{i,t-1}/L_{it}$.
Agents form rational expectations $\mathbb{E}_t[\phi_{i,t+H}]$ at horizon $H$.
\end{assumption}

\begin{assumption}[Effective Retirement Duration]\label{ass:A3}
Old-age length, normalised by working-life length, equals $\omega_{it}$.
Agents form rational expectations $\mathbb{E}_t[\omega_{i,t+H}]$ at horizon $H$.
\end{assumption}

\begin{assumption}[Two-Country World]\label{ass:A4}
The world consists of two countries $i\in\{H,F\}$ (the analysis extends to $N$
countries with no change in structure). Asset markets are perfectly integrated:
all assets earn the world rate $r_t^*$.
\end{assumption}

\begin{assumption}[Technology]\label{ass:A5}
Output of country $i$ at $t$ is Cobb--Douglas: $Y_{it}=K_{it}^{\alpha}
(A_{it}L_{it})^{1-\alpha}$, with $\alpha\in(0,1)$ and depreciation $\delta$.
Define $g_{it}^A\equiv \Delta \ln A_{it}$, $g_{it}^L\equiv \Delta\ln L_{it}$,
$g_{it}^Y\equiv g_{it}^A+g_{it}^L$; on the BGP these converge to country-specific
constants $g_i^A, g_i^L, g_i^Y$.
\end{assumption}

\begin{assumption}[Preferences]\label{ass:A6}
Lifetime utility is CRRA:
$U=u(c_{it}^y)+\omega_{it}\beta\,\mathbb{E}_t u(c_{i,t+1}^o)$ with
$u(c)=c^{1-1/\sigma}/(1-1/\sigma)$, $\sigma>0,\sigma\ne1$.
\end{assumption}

\begin{assumption}[No Bequests]\label{ass:A7}
Agents derive no utility from end-of-life wealth and face no bequest motive.
Equivalently, the terminal condition on old-age financial wealth is
$a_{i,t+1}^{o\,\text{end}}=0$. This assumption is invoked explicitly in the
household derivation (\cref{sec:theory:households}), in the aggregation of
old-age dissaving (\cref{sec:theory:agg-saving}), and in the partial-Ricardian
argument (\cref{rem:ricardian}).
\end{assumption}

\begin{assumption}[Government Budget]\label{ass:A8}
The government of country $i$ collects payroll taxes $\tau_{it}$ from the
young, pays pension transfers $tr_{it}$ to the old, and issues debt $D_{it}^g$.
The flow budget per efficiency worker, defining $n_{it}\equiv\Delta\ln L_{it}$,
is
\begin{equation}
d_{i,t+1}^g(1+g_{it}^Y)=d_{it}^g(1+r_t^*)+\phi_{it}\,tr_{it}-\tau_{it},
\end{equation}
where $d_{it}^g\equiv D_{it}^g/(A_{it}L_{it})$.
\end{assumption}

\begin{assumption}[Stability of the World Capital Market]\label{ass:A9}
The excess world saving function $ES(r_t^*;\mathbf Z_t)$ (defined in
\cref{sec:theory:clearing}) satisfies $\partial ES/\partial r_t^*>0$ in a
neighbourhood of equilibrium. Equivalently, the net supply of world loanable
funds is locally increasing in the world rate.
\end{assumption}

\Cref{ass:A9} is the regularity condition that licenses the IFT used in
\cref{prop:cs-rstar}. It holds when investment is sufficiently elastic in
$r^*$ relative to the (potentially offsetting income- and
substitution-effect-driven) response of saving; it is standard in
OLG comparative-statics arguments \citep{Diamond1965}.

\subsection{Households: Life-Cycle Optimisation}\label{sec:theory:households}

A young agent in country $i$ at $t$ chooses $(c_{it}^y,s_{it})$ to maximise
\[
U_{it}=u(c_{it}^y)+\omega_{it}\beta\,\mathbb{E}_t u(c_{i,t+1}^o)
\]
subject to the young- and old-age budget constraints
\begin{align}
c_{it}^y+s_{it}&=w_{it}-\tau_{it}, \\
\omega_{it}\,c_{i,t+1}^o&=(1+r_{t+1}^*)\,s_{it}+tr_{i,t+1},
\end{align}
and the no-bequest terminal condition (\cref{ass:A7}): all financial wealth
$(1+r_{t+1}^*)s_{it}+tr_{i,t+1}$ is consumed by end of the old period.
Without \cref{ass:A7}, the old's consumption could fall short of accumulated
wealth and a bequest term would appear in the aggregate-saving equation; we
return to this when aggregating in \cref{sec:theory:agg-saving}.

Substituting the old-age budget into the lifetime constraint:
\[
\Omega_{it}\equiv (w_{it}-\tau_{it})+\frac{tr_{i,t+1}}{1+r_{t+1}^*}.
\]
The CRRA Euler equation gives the closed form
\begin{equation}
c_{it}^{y*}=\mu(r_{t+1}^*,\omega_{it})\,\Omega_{it},\qquad
\mu\equiv\frac{1}{1+\omega_{it}\beta^{1/\sigma}(1+r_{t+1}^*)^{(1-\sigma)/\sigma}}.
\label{eq:mu}
\end{equation}
Individual saving is then
\begin{equation}
s_{it}^*=(w_{it}-\tau_{it})-\mu(r_{t+1}^*,\omega_{it})\,\Omega_{it}.
\label{eq:savind}
\end{equation}

\begin{proposition}[Individual Saving: Signed Comparative Statics]\label{prop:indiv-saving}
At an interior optimum:
\begin{enumerate}
\item $\partial s_{it}^*/\partial(w_{it}-\tau_{it})\in(0,1)$;
\item $\partial s_{it}^*/\partial \omega_{it}>0$;
\item $\partial s_{it}^*/\partial \,tr_{i,t+1}=-\mu/(1+r_{t+1}^*)<0$,
holding $(r_{t+1}^*,\omega_{it})$ fixed;
\item $\partial s_{it}^*/\partial r_{t+1}^*=-(\partial\mu/\partial r_{t+1}^*)\Omega_{it}+\mu\,tr_{i,t+1}/(1+r_{t+1}^*)^2$, with
\[
\frac{\partial\mu}{\partial r_{t+1}^*}=-\omega_{it}\beta^{1/\sigma}\frac{1-\sigma}{\sigma}(1+r_{t+1}^*)^{\frac{1-2\sigma}{\sigma}}\Big/[1+\omega_{it}\beta^{1/\sigma}(1+r_{t+1}^*)^{(1-\sigma)/\sigma}]^2,
\]
so the sign is unambiguously positive for $\sigma\le 1$ (substitution effect
dominates) and ambiguous for $\sigma>1$.
\end{enumerate}
\end{proposition}

\begin{proof}
Parts (1)--(3) follow by direct differentiation of \eqref{eq:savind}.
Part (4) writes out the derivative explicitly; the sign of $\partial\mu/\partial r^*$ is the sign of $\sigma-1$. The second term is unambiguously positive (the present value of transfers falls in $r^*$, raising desired saving by the young). The first term is non-negative iff $\sigma\le 1$.
\end{proof}

\subsection{Government Sector and National Saving; Partial Ricardian Offset}
\label{sec:theory:gov}

National saving decomposes as
$S_{it}^{nat}=S_{it}^{priv}+S_{it}^g$ where $S_{it}^g\equiv \tau_{it}L_{it}-tr_{it}L_{i,t-1}-r_t^*D_{i,t-1}^g$ is the structural primary balance minus interest service.

\begin{remark}[Partial Ricardian Offset]\label{rem:ricardian}
Under \crefrange{ass:A2}{ass:A8} (two-period OLG, no bequests), private
saving rises by \emph{less} than one-for-one in response to public debt:
$\partial S^{priv}/\partial D^g\in (-1,0)$, hence
$\partial S^{nat}/\partial D^g=\partial S^{priv}/\partial D^g-1\in(-1,0)$.
\end{remark}

\begin{proof}[Proof sketch]
In a Diamond two-period OLG with \cref{ass:A7} (no bequests), the young do
not internalise the old's intertemporal budget: government bonds issued today
finance transfers to today's old, who are not the young's heirs. Hence the
young treat bonds as net wealth and do not increase saving by the full present
value of future taxes \citep{Diamond1965,Blanchard1985,Buiter1988}. Formally,
differentiating $s_{it}^*=(w_{it}-\tau_{it})-\mu\Omega_{it}$ with respect to a
debt-financed cut in $\tau_{it}$ (with a future tax increase falling on the
young's grandchildren, not on themselves), the young's saving response is
strictly less than the present value of the future tax burden. Thus
$\gamma_D\equiv (r^*-g_i^Y)^{-1}\,\partial S^{nat}/\partial d^g\in
(-(r^*-g_i^Y)^{-1},0)$, so $\gamma_D<0$ but bounded above zero by partial
offset.
\end{proof}

\subsection{Aggregate Saving}\label{sec:theory:agg-saving}

We now aggregate over young and old to obtain private saving per young worker.
Throughout, $s_{it}^*$ denotes the \emph{individual} saving of a single young
agent of generation $t$ (not a per-worker aggregate). All young agents in
country $i$ are ex-ante identical, so per-young-worker young saving equals
$s_{it}^*$.

\paragraph{Young.} There are $L_{it}$ identical young, each saving $s_{it}^*$.
Per young worker their saving is exactly $s_{it}^*$.

\paragraph{Old.} There are $L_{i,t-1}$ old in period $t$. By \cref{ass:A7},
each old agent entered period $t$ with financial wealth
$(1+r_t^*)s_{i,t-1}^*$, also receives transfer $tr_{it}$, and consumes
$(1+r_t^*)s_{i,t-1}^*+tr_{it}$ entirely. The old's \emph{gross dissaving} of
\emph{financial} wealth (the only part that contributes to private saving;
transfers are a redistribution accounted for in $S_{it}^g$) is
$L_{i,t-1}(1+r_t^*)s_{i,t-1}^*$.

\paragraph{Per young worker.} Aggregate private saving per young worker is the
young's saving minus the old's gross financial dissaving, scaled by $L_{i,t-1}/L_{it}=\phi_{it}$:
\begin{equation}
\frac{S_{it}^{priv}}{L_{it}}=s_{it}^*-\phi_{it}\,(1+r_t^*)\,s_{i,t-1}^*.
\label{eq:agg-priv-saving}
\end{equation}
The transfer term $tr_{it}$ does \emph{not} appear in \eqref{eq:agg-priv-saving}
because, on the consolidated private--government accounts, transfers are an
intragenerational redistribution: they reduce $S_{it}^g$ by $\phi_{it}\,tr_{it}$
and raise old consumption by the same amount, leaving \emph{national} saving
unaffected. This treatment is gross of transfer receipts by the old; the
government saving block in \cref{sec:theory:gov} accounts for the same flows
on the public side. Without \cref{ass:A7}, the old would retain a bequest
$\xi_{it}>0$, and \eqref{eq:agg-priv-saving} would gain a term
$\phi_{it}\xi_{it}$.

\paragraph{Reduced-form saving function.} Writing $s_{it}^*$ as a function of
its determinants from \cref{prop:indiv-saving} and substituting the competitive
wage, define the structural per-worker saving function
\begin{equation}
s_{it}^Y\equiv\mathcal{S}\bigl(y_{it},\phi_{it},\mathbb{E}_t[\phi_{i,t+H}],\mathbb{E}_t[\omega_{i,t+H}],d_{it}^g,r_t^*\bigr),
\label{eq:S-function}
\end{equation}
where $y_{it}\equiv \ln Y_{it}/(A_{it}L_{it})$ is log output per efficiency
worker (Modigliani--Brumberg life-cycle channel; \citealp{ModiglianiBrumberg1954,Feldstein1980}).
The signed comparative statics of $\mathcal{S}$ are inherited from
\cref{prop:indiv-saving} and the aggregation in \eqref{eq:agg-priv-saving}:
$\partial\mathcal{S}/\partial y>0$, $\partial\mathcal{S}/\partial\phi<0$
(current retirees dissave), $\partial\mathcal{S}/\partial \mathbb{E}_t[\phi_{i,t+H}]>0$
(precautionary), $\partial\mathcal{S}/\partial \mathbb{E}_t[\omega_{i,t+H}]>0$
(retirement-duration), $\partial\mathcal{S}/\partial d^g<0$ (partial Ricardian
offset, \cref{rem:ricardian}).

\subsection{Production, Wages, Productivity}\label{sec:theory:production}

With Cobb--Douglas technology (\cref{ass:A5}) and competitive factor markets,
the marginal product of capital equals $r_t^*+\delta$, which yields a
capital--labour ratio (per efficiency worker) $\hat k(r_t^*)=[\alpha/(r_t^*+\delta)]^{1/(1-\alpha)}$
and competitive wage
\begin{equation}
w_{it}=(1-\alpha)\left(\frac{\alpha}{r_t^*+\delta}\right)^{\alpha/(1-\alpha)} A_{it}.
\label{eq:wage}
\end{equation}
Hence wages depend on the world rate $r_t^*$ (common across countries) and on
country-specific TFP $A_{it}=\bar A_t\exp(\tilde a_{it})$, where $\tilde a_{it}$
is the log deviation from the world mean. Output per efficiency worker on the
BGP is therefore a function of $r^*$ and $\tilde a_i$ alone.

\subsection{Firms and Investment}\label{sec:theory:firms}

Investment is governed by Tobin's $q$ with quadratic adjustment costs
\citep{Hayashi1982}. Per efficiency worker, the firm chooses $\iota_{it}$ to
maximise
\[
\sum_{s\ge 0}\Bigl[\prod_{j=0}^{s-1}(1+r_{t+j}^*)^{-1}\Bigr]\Bigl[F(k_{i,t+s},1)-\iota_{i,t+s}-\frac{\chi}{2}(\iota_{i,t+s}-\delta\,k_{i,t+s})^2/k_{i,t+s}\Bigr],
\]
yielding the Euler condition $q_{it}=1+\chi(\iota_{it}/k_{it}-\delta)$. Along
the BGP, $\iota_{it}/k_{it}\to \delta+g_i^Y/(1+\alpha\chi)$ and the structural
per-worker investment function is
\begin{equation}
\iota_{it}\equiv\mathcal{I}\bigl(\tilde a_{it},\mathbb{E}_t[\phi_{i,t+H}],\mathbb{E}_t[\omega_{i,t+H}],r_t^*\bigr),
\label{eq:I-function}
\end{equation}
with $\partial\mathcal{I}/\partial r^*<0$, $\partial\mathcal{I}/\partial\tilde a>0$
(higher MPK raises investment), and $\partial\mathcal{I}/\partial\mathbb{E}_t[\phi_{i,t+H}]<0$
(anticipated lower future labour input lowers desired capital).

\subsection{Global Capital-Market Clearing}\label{sec:theory:clearing}

World saving must equal world investment plus the world stock of public debt
financed externally:
\begin{equation}
\sum_i \eta_{it}\mathcal{S}(\cdot)=\sum_i \eta_{it}\mathcal{I}(\cdot)+\sum_i \eta_{it}\,d_{it}^g,
\label{eq:world-clearing}
\end{equation}
where $\eta_{it}\equiv A_{it}L_{it}/\sum_j A_{jt}L_{jt}$ are GDP weights. Define
\[
ES(r_t^*;\mathbf Z_t)\equiv \sum_i\eta_{it}\mathcal{S}(\cdot)-\sum_i\eta_{it}\mathcal{I}(\cdot)-\sum_i\eta_{it}\,d_{it}^g,
\]
where $\mathbf Z_t$ collects all country-time fundamentals other than $r_t^*$.
Equilibrium $r_t^*$ solves $ES(r_t^*;\mathbf Z_t)=0$.

\begin{proposition}[Comparative Statics of $r_t^*$]\label{prop:cs-rstar}
Under \cref{ass:A9} ($\partial ES/\partial r^*>0$), the IFT applied to
$ES(r_t^*;\mathbf Z_t)=0$ yields, for any determinant $z\in\mathbf Z_t$,
\[
\frac{\partial r_t^*}{\partial z}=-\frac{\partial ES/\partial z}{\partial ES/\partial r^*}.
\]
The denominator is positive by \cref{ass:A9}; signs of the four claims below
follow from the explicit numerators.
\begin{enumerate}
\item[(i)] \emph{Anticipated global ageing} (conditional). With
$\partial \mathcal{S}/\partial \mathbb{E}_t[\phi^e]>0$ and
$\partial \mathcal{I}/\partial \mathbb{E}_t[\phi^e]<0$, the numerator is
$\partial ES/\partial \mathbb{E}_t[\phi^e]>0$ \emph{provided the saving
response dominates the investment response in magnitude}. Under that
condition, $\partial r^*/\partial \mathbb{E}_t[\phi^e]<0$.
\item[(ii)] \emph{Global public debt} (unconditional under OLG). Each unit of
$d_i^g$ enters $ES$ with coefficient $-1$ directly, and the indirect effect on
$\mathcal{S}$ is bounded above by zero by partial Ricardian offset
(\cref{rem:ricardian}). Hence $\partial ES/\partial d^g<0$ \emph{net of
private-saving response}, so $\partial r^*/\partial d^g>0$.
\item[(iii)] \emph{Global TFP} (ambiguous). $\partial \mathcal{S}/\partial\tilde a>0$
(higher wage raises young saving) and $\partial\mathcal{I}/\partial\tilde a>0$
(higher MPK raises investment); the sign of $\partial ES/\partial \tilde a$ is
indeterminate.
\item[(iv)] \emph{Current dependency} (ambiguous). $\partial\mathcal{S}/\partial\phi<0$
(retiree dissaving) and $\partial\mathcal{I}/\partial\phi$ is itself
indeterminate; the sign of $\partial r^*/\partial\phi$ is therefore not pinned
down by theory.
\end{enumerate}
\end{proposition}

\begin{proof}
The denominator is positive by \cref{ass:A9}. For each claim, we compute the
numerator $\partial ES/\partial z$ by differentiating \eqref{eq:world-clearing}
and substituting the comparative statics of $\mathcal{S}$
(\cref{sec:theory:agg-saving}) and $\mathcal{I}$ (\cref{sec:theory:firms}).
\emph{(i)} $\partial ES/\partial \mathbb{E}_t[\phi^e]=\sum_i\eta_i[\partial\mathcal{S}/\partial \mathbb{E}_t[\phi^e]-\partial\mathcal{I}/\partial \mathbb{E}_t[\phi^e]]$.
Both bracketed terms are positive; the sum is positive when the saving channel
dominates investment. \emph{(ii)} $\partial ES/\partial d^g=\sum_i\eta_i[\partial\mathcal{S}/\partial d^g-1]<0$
because $\partial\mathcal{S}/\partial d^g\in (-1,0)$ by \cref{rem:ricardian}.
\emph{(iii), (iv)} Both signed components enter with opposite signs in
$ES$; signs are not pinned down without further restrictions.
\end{proof}

\begin{corollary}[GE Spillover]\label{cor:ge-spillover}
A shock to country-$F$ fundamentals affects country-$H$ NIIP only through
$r_t^*$. Because $r_t^*$ is common across countries at any date, the GE
spillover is absorbed by time fixed effects in a panel regression. This is
formalised in \cref{sec:theory:linearization}, Step (d).
\end{corollary}

\subsection{NIIP Law of Motion}\label{sec:theory:niip-lom}

Let $B_{it}$ denote the level NIIP and $b_{it}\equiv B_{it}/(A_{it}L_{it})$
the per-efficiency-worker NIIP. The balance-of-payments identity, in
per-efficiency-worker form, is
\begin{equation}
b_{i,t+1}=\frac{1+r_t^*}{1+g_{it}^Y}\,b_{it}+s_{it}^Y-\iota_{it}+\text{residuals},
\label{eq:niip-lom-discrete}
\end{equation}
where residuals collect valuation effects and capital transfers
\citep{ObstfeldRogoff1995}. Linearising around small rates and writing
$R_{it}\equiv (1+r_t^*)/(1+g_{it}^Y)\approx 1+r_t^*-g_{it}^Y$:
\begin{equation}
\Delta b_{i,t+1}\approx (r_t^*-g_{it}^Y)\,b_{it}+ca_{it}^{NR}+\text{residuals},
\qquad ca_{it}^{NR}\equiv s_{it}^Y-\iota_{it}.
\label{eq:niip-lom-cts}
\end{equation}
We use \eqref{eq:niip-lom-discrete} in derivations and \eqref{eq:niip-lom-cts}
for intuition.

\subsection{Intertemporal Solvency}\label{sec:theory:solvency}

Iterating \eqref{eq:niip-lom-discrete} forward and imposing the no-Ponzi
condition $\lim_{S\to\infty}\bigl[\prod_{j=0}^{S-1}R_{i,t+j}\bigr]^{-1}b_{i,t+S}=0$
gives the country-$i$ intertemporal budget constraint
\begin{equation}
b_{it}=\sum_{s=1}^{\infty}\Bigl[\prod_{j=0}^{s-1}R_{i,t+j}\Bigr]^{-1}\bigl(\iota_{i,t+s}-s_{i,t+s}^Y\bigr).
\label{eq:ibc}
\end{equation}
\Cref{eq:ibc} states that the current NIIP equals the discounted sum of future
non-interest current-account deficits \citep{Sachs1981,ObstfeldRogoff1995}.

\subsection{From IBC to the Structural Long-Run NIIP Equation}
\label{sec:theory:linearization}

This subsection contains the central derivation of the paper. It was missing
from the original draft and is reconstructed here in four steps: (a) the IBC
collapses on the BGP to a perpetuity; (b) structural saving and investment are
substituted; (c) a first-order Taylor expansion produces signed coefficients;
(d) the deviation of $r_t^*$ from its BGP value is shown to be common across
countries at any date $t$, so time fixed effects $\lambda_t$ absorb the GE
channel.

\subsubsection*{Step (a): BGP collapse}

On the BGP (\cref{ass:A2,ass:A5}): $r_t^*\to r^*$, $g_{it}^Y\to g_i^Y$, both
constant; structural saving and investment converge to country-specific
constants $\bar s_i,\bar\iota_i$. The discount factor becomes
$\bar R_i\equiv(1+r^*)/(1+g_i^Y)$, constant. The no-Ponzi condition holds
under dynamic efficiency ($r^*>g_i^Y$, equivalent to $\bar R_i>1$). The IBC
\eqref{eq:ibc} then collapses to a geometric series:
\begin{align}
b_i^{LR}&=\sum_{s=1}^{\infty}\bar R_i^{-s}(\bar s_i-\bar\iota_i)\notag\\
&=\frac{\bar R_i^{-1}}{1-\bar R_i^{-1}}(\bar s_i-\bar\iota_i)\notag\\
&=\frac{1}{\bar R_i-1}\,\bar{ca}_i^{NR}\notag\\
&=\frac{\bar{ca}_i^{NR}}{r^*-g_i^Y},
\label{eq:bgp-niip}
\end{align}
where we have used $\bar R_i-1\approx r^*-g_i^Y$ and defined
$\bar{ca}_i^{NR}\equiv \bar s_i-\bar\iota_i$. \Cref{eq:bgp-niip} is the
\emph{perpetuity identity}: the long-run NIIP equals the present value of a
perpetual stream of non-interest current-account surpluses.

\subsubsection*{Step (b): Structural substitution}

Substituting the reduced-form structural saving and investment functions
\eqref{eq:S-function}--\eqref{eq:I-function} into \eqref{eq:bgp-niip}:
\begin{equation}
b_i^{LR}=\mathcal{B}(r^*,g_i^Y,y_i,\phi_i,\mathbb{E}[\phi_{i+H}],\mathbb{E}[\omega_{i+H}],d_i^g,\tilde a_i),
\label{eq:structural-niip}
\end{equation}
where
\[
\mathcal{B}(\cdot)\equiv \frac{\mathcal{S}(\cdot)-\mathcal{I}(\cdot)}{r^*-g_i^Y}.
\]
At this stage, $r^*$ is treated as a parameter; in Step (d) we will recognise
that the world-clearing condition makes $r^*$ a function of all countries'
fundamentals.

\subsubsection*{Step (c): First-order Taylor expansion around country-specific BGP}

Expand $\mathcal{B}(\cdot)$ around the country-specific BGP reference point
$\bar{\mathbf x}_i=(\bar y_i,\bar\phi_i,\overline{\mathbb{E}}[\phi_{i+H}],\overline{\mathbb{E}}[\omega_{i+H}],\bar d_i^g,\bar{\tilde a}_i)$
and the world reference rate $\bar r^*$. Writing $x_{it}^k$ for the $k$th
country-time fundamental,
\begin{equation}
b_{it}^{LR}\approx \mathcal{B}(\bar{\cdot})+\sum_k \gamma_k^i\,(x_{it}^k-\bar x_i^k)+\frac{\partial\mathcal{B}}{\partial r^*}\,(r_t^*-\bar r^*),
\label{eq:taylor}
\end{equation}
where the structural coefficients are
\begin{equation}
\gamma_k^i=\frac{1}{r^*-g_i^Y}\left.\frac{\partial(\mathcal{S}-\mathcal{I})}{\partial x_k}\right|_{\bar{\mathbf x}_i}.
\label{eq:gamma}
\end{equation}
The scalar $1/(r^*-g_i^Y)>0$ under dynamic efficiency, so $\gamma_k^i$ inherits
the sign of $\partial(\mathcal{S}-\mathcal{I})/\partial x_k$ at the BGP.
Reading off from the comparative statics of $\mathcal{S}$
(\cref{sec:theory:agg-saving}) and $\mathcal{I}$ (\cref{sec:theory:firms}):
\Cref{tab:signs} lists the seven structural sign predictions, each derived,
not asserted.

\begin{table}[t]
\centering
\caption{Sign Predictions for Structural Coefficients}\label{tab:signs}
\begin{threeparttable}
\begin{tabular}{lll}
\toprule
Coefficient & Channel & Predicted sign \\
\midrule
$\gamma_y$ on $y_{it}$ & Life-cycle: higher income raises young saving & $>0$ \\
$\gamma_\phi$ on $\phi_{it}$ & Current retiree dissaving dominates & $<0$ \\
$\gamma_{\phi^e}$ on $\mathbb{E}_t[\phi_{i,t+H}]$ & Precautionary saving; lower future labour & $>0$ \\
$\gamma_{\omega^e}$ on $\mathbb{E}_t[\omega_{i,t+H}]$ & Retirement-duration motive & $>0$ \\
$\gamma_{TFP}$ on $\tilde a_{it}$ & Saving and investment both rise; decomposed in App. & ambiguous \\
$\gamma_D$ on $d_{it}^g$ & Partial Ricardian offset (\cref{rem:ricardian}) & $\in(-(r^*-g^Y)^{-1},0)$ \\
$\partial\mathcal{B}/\partial r^*$ on $r_t^*$ & GE channel: absorbed by $\lambda_t$ in Step (d) & $-$ \\
\bottomrule
\end{tabular}
\begin{tablenotes}\footnotesize
\item \textit{Notes.} Each sign is derived from the explicit form of
$\mathcal{S}$ and $\mathcal{I}$, not asserted. Ambiguous TFP is decomposed
into a saving, investment, and GE component in the Technical Appendix
(\cref{sec:theory:appendix}).
\end{tablenotes}
\end{threeparttable}
\end{table}

\subsubsection*{Step (d): $r_t^*$ deviation is common at date $t\Rightarrow \lambda_t$ absorbs the GE channel}

From the world-clearing condition \eqref{eq:world-clearing}, applying the IFT
of \cref{prop:cs-rstar} along the BGP,
\[
r_t^*-\bar r^*=\sum_i \eta_i\,\boldsymbol\psi_i^{\prime}\,(\mathbf x_{it}-\bar{\mathbf x}_i),
\]
where $\boldsymbol\psi_i$ collects the IFT-derived GE weights. The key
observation is that the right-hand side depends only on country-time
fundamentals weighted into a \emph{world} aggregate; once aggregated, the
result is the \emph{same number for every country $i$ at any given date $t$}.
That is,
\[
r_t^*-\bar r^* \;\text{ is a function of }\, (\mathbf x_{1t},\ldots,\mathbf x_{Nt}) \;\text{ but not of }\, i.
\]
Hence the term $(\partial\mathcal{B}/\partial r^*)(r_t^*-\bar r^*)$ in
\eqref{eq:taylor} is the \emph{same} at any date $t$ across all countries.

In a panel regression with time fixed effects $\lambda_t$, any term that
varies only across $t$ (and not across $i$) is exactly absorbed by $\lambda_t$.
Substituting into \eqref{eq:taylor}, taking $\gamma_k^i\approx \gamma_k$ as a
common slope at the world average BGP (with country-specific deviations
absorbed in $\mu_i$), we obtain the structural regression equation
\begin{empheq}[box=\fbox]{equation}
b_{it}^{LR}=\gamma_y\,y_{it}+\gamma_\phi\,\phi_{it}+\gamma_{\phi^e}\,\mathbb{E}_t[\phi_{i,t+H}]+\gamma_{\omega^e}\,\mathbb{E}_t[\omega_{i,t+H}]+\gamma_{TFP}\,\tilde a_{it}+\gamma_D\,d_{it}^g+\mu_i+\lambda_t+\varepsilon_{it},
\label{eq:structural-reg}
\end{empheq}
where:
\begin{itemize}[noitemsep]
\item $\mu_i\equiv \mathcal{B}(\bar{\cdot})-\sum_k \gamma_k\bar x_i^k$ absorbs time-invariant country-specific BGP deviations (including country-specific $g_i^Y$ scaling);
\item $\lambda_t\equiv (\partial\mathcal{B}/\partial r^*)(r_t^*-\bar r^*)$ absorbs the world interest rate deviation and any other globally common shock;
\item $\varepsilon_{it}$ captures higher-order linearisation error, country-specific slope deviations from $\gamma_k^i\approx\gamma_k$, and measurement error.
\end{itemize}

\Cref{eq:structural-reg} is the structural counterpart of the empirical
specification estimated in \cref{sec:results} (forward reference to the
empirical section). It is the empirical promise of the model: each
$\gamma_k$ is interpretable as a partial derivative of $\mathcal{S}-\mathcal{I}$
scaled by $(r^*-g^Y)^{-1}$, evaluated at the BGP. The within-time
interpretation of $\hat\gamma_k$ in panel estimation is exactly licensed by
Step (d).

\subsection{Error-Correction Representation}\label{sec:theory:ecm}

If $b_{it}$ and the regressors $(y_{it},\phi_{it},\ldots)$ are $I(1)$ and
cointegrate via \eqref{eq:structural-reg}, the Granger representation theorem
\citep{EngleGranger1987} delivers the error-correction (ECM) form
\begin{equation}
\Delta b_{it}=\rho_i\bigl[b_{i,t-1}-(\gamma_y y_{i,t-1}+\cdots+\gamma_D d_{i,t-1}^g+\mu_i+\lambda_t)\bigr]+\boldsymbol\beta_i^{\prime}\Delta\mathbf x_{it}+\zeta_{it},
\label{eq:ecm}
\end{equation}
where:
\begin{itemize}[noitemsep]
\item $\rho_i\in(-1,0)$ is the country-$i$ \emph{adjustment speed} (this is
the parameter previously denoted $\alpha_i$ and renamed to free $\alpha$ for
the capital share);
\item $\boldsymbol\beta_i$ collects short-run coefficients on first differences
of the regressors;
\item $\zeta_{it}$ is the short-run innovation.
\end{itemize}
The ECM additionally permits short-run regressors that are mean-reverting
($I(0)$) and therefore enter \eqref{eq:ecm} but not the long-run
\eqref{eq:structural-reg}. We allow two such regressors:
$pb_{it}^{str}\equiv \tau_{it}-\phi_{it}tr_{it}$, the structural primary
balance per efficiency worker (which is a flow, not a stock); and
$\Delta\ln q_{it}$, the change in the real effective exchange rate (which is
mean-reverting under PPP). Both are absent from \eqref{eq:structural-reg}
because the long-run NIIP depends on the \emph{stock} of public debt $d^g$
(not the primary-balance flow) and is invariant to nominal exchange rates.

\subsection{Summary: Seven Structural Predictions}\label{sec:theory:summary}

The model delivers seven testable structural predictions:

\begin{enumerate}[noitemsep]
\item $\gamma_y>0$: higher income per efficiency worker raises long-run NIIP.
\item $\gamma_\phi<0$: a higher current dependency ratio lowers long-run NIIP.
\item $\gamma_{\phi^e}>0$: anticipated future ageing raises long-run NIIP through precautionary saving by the young.
\item $\gamma_{\omega^e}>0$: anticipated longer retirement durations raise long-run NIIP.
\item $\gamma_{TFP}$: \emph{a priori} ambiguous; decomposed in \cref{sec:theory:appendix} into saving, investment, and GE channels.
\item $\gamma_D\in(-(r^*-g^Y)^{-1},0)$: public debt lowers long-run NIIP, but \emph{less than one-for-one} (partial Ricardian offset; \cref{rem:ricardian}).
\item Time fixed effects $\lambda_t$ absorb the world real interest rate
$r_t^*$ and all global common shocks (\cref{cor:ge-spillover}, Step (d) of
\cref{sec:theory:linearization}).
\end{enumerate}

These seven predictions map one-to-one onto the panel coefficients estimated
in \cref{sec:results}.

\subsection{Technical Appendix: TFP Decomposition}\label{sec:theory:appendix}

The ambiguous sign of $\gamma_{TFP}$ in \cref{tab:signs} reflects three offsetting channels. Substituting \eqref{eq:S-function}--\eqref{eq:I-function} into the
expression $\gamma_{TFP}^i=(r^*-g_i^Y)^{-1}\partial(\mathcal{S}-\mathcal{I})/\partial\tilde a$, and recognising that $r^*$ itself responds to global TFP through \cref{prop:cs-rstar}(iii),
\begin{equation}
\frac{db_i^{LR}}{d\tilde a_i}=\underbrace{\frac{1}{r^*-g_i^Y}\frac{\partial\mathcal{S}}{\partial\tilde a_i}}_{\text{(+) saving channel}}-\underbrace{\frac{1}{r^*-g_i^Y}\frac{\partial\mathcal{I}}{\partial\tilde a_i}}_{\text{(+) investment channel (\(-\) on \(b^{LR}\))}}+\underbrace{\frac{\partial\mathcal{B}}{\partial r^*}\frac{\partial r^*}{\partial\bar a}\eta_i}_{\text{(?) GE channel}}.
\label{eq:tfp-decomp}
\end{equation}
The first term raises $b^{LR}$ (higher wages, more young saving); the second
lowers it (higher MPK, more investment demand crowds out NFA); the third is
the GE feedback through the world rate, weighted by country $i$'s world share
$\eta_i$ and absorbed by $\lambda_t$ in the panel. Empirically, the second
channel is expected to dominate in catching-up economies; the first dominates
in mature, saving-rich economies.

% End of theoretical_model.tex
